% AY2017 材料力学 〔I〕（転記）
% 出典: 04_材料力学/original/04_材料力学/2017/2017za.pdf
% 要約: 丸棒\textcircled{1}(断面積2A, 長さL1, ヤング率2E)と2本の丸棒\textcircled{2}(断面積A, 長さL2, ヤング率E)を
%       剛体板に接合。丸棒\textcircled{1}が\textcircled{2}よりわずかに短いため, 高さL3(>L2-L1)の剛体円盤(断面積2A)を
%       はめ込んだ。棒・剛体板の自重は無視, 剛体板・剛体円盤は変形しない。
%       (1) σ1,σ2, (2) λ1,λ2, (3) 上面剛体板中央に力を加え丸棒\textcircled{2}を元の長さL2に戻す外力P。

\section*{〔I〕}

図1に示すように, 丸棒\textcircled{1}と 2 本の丸棒\textcircled{2}を剛体板に接合した. 丸棒\textcircled{1}が左右の丸棒\textcircled{2}に比べて
わずかに短いため, 高さ $L_3\,(>L_2-L_1)$ の剛体円盤（断面積 $2A$）をはめ込んだ.
次の問い (1)\,$\sim$\,(3) に答えよ. ただし, 棒や剛体板の自重は無視するものとし,
剛体板, 剛体円盤は変形しないものとする.

\begin{center}
\begin{tikzpicture}[>=Latex,scale=1.0]
  % 下面: 地面 + 下剛体板
  \draw[thick] (-0.4,0) -- (5.4,0);
  \foreach \x in {-0.3,0.0,...,5.3} {\draw (\x,0) -- (\x-0.22,-0.22);}
  \fill[gray!20] (0.2,0) rectangle (4.8,0.4); \draw (0.2,0) rectangle (4.8,0.4);
  % 上剛体板
  \fill[gray!20] (0.2,4.2) rectangle (4.8,4.6); \draw (0.2,4.2) rectangle (4.8,4.6);
  % 丸棒\textcircled{2}（左, 全長 L2: y0.4..4.2）
  \draw (0.7,0.4) rectangle (1.2,4.2);
  % 丸棒\textcircled{2}（右）
  \draw (3.8,0.4) rectangle (4.3,4.2);
  % 丸棒\textcircled{1}（中央, 短い L1: y0.4..3.2）
  \draw (2.1,0.4) rectangle (2.9,3.2);
  % 剛体円盤（中央上, 高さL3: y3.2..4.2）
  \fill[gray!35] (2.1,3.2) rectangle (2.9,4.2); \draw (2.1,3.2) rectangle (2.9,4.2);
  % ラベル
  \node at (0.95,3.6) {\textcircled{2}}; \node at (4.05,3.6) {\textcircled{2}}; \node at (2.5,1.7) {\textcircled{1}};
  % 寸法 L2（左外）, L1（中央内）
  \draw[<->] (-0.1,0.4) -- (-0.1,4.2) node[midway,fill=white]{$L_2$};
  \draw[<->] (1.6,0.4) -- (1.6,3.2) node[midway,fill=white]{$L_1$};
  % 剛体板ラベル
  \node[left] at (0.2,4.4) {\footnotesize 剛体板};
  \node[left] at (0.2,0.2) {\footnotesize 剛体板};
\end{tikzpicture}
\qquad
\begin{tabular}{@{}l@{}}
  丸棒\textcircled{1}\\ 断面積 $2A$, 長さ $L_1$, ヤング率 $2E$\\[6pt]
  丸棒\textcircled{2}\\ 断面積 $A$, 長さ $L_2$, ヤング率 $E$
\end{tabular}

\vspace{1mm}
図1
\end{center}

\begin{enumerate}[label=(\arabic*)]
  \item 丸棒\textcircled{1}, \textcircled{2}に生じる応力 $\sigma_1$, $\sigma_2$ を求めよ.

  \item 丸棒\textcircled{1}, \textcircled{2}の変位 $\lambda_1$, $\lambda_2$ を求めよ.

  \item 上面の剛体板中央に力を加え, 丸棒\textcircled{2}を元の長さ $L_2$ まで戻すために必要な
        外力 $P$ を求めよ.
\end{enumerate}
